% Agent-authored source transcription and faithful translation, verified against PDF 25-28.
% The sample starts at Section 2.4; it excludes the preceding Section 2.3 and later Box 2.1.
% block: chain-maps-heading page:25 type:section
\subsection{\BT{Chain maps}{链映射}}\label{sec:chain-maps}

% block: chain-map-definition page:25 type:ordinary-paragraph
\BT{A \emph{chain map} from a chain complex $(V,\partial)$ to a second chain complex $(W,\xi)$ is a sequence of linear maps $f_k:V_k\to W_k$ such that the following diagram commutes:}{从链复形 $(V,\partial)$ 到另一个链复形 $(W,\xi)$ 的\emph{链映射}是一列线性映射 $f_k:V_k\to W_k$, 使得下图交换:}

% block: chain-map-diagram page:25 type:equation
\[
\begin{tikzcd}[column sep=1.2em,row sep=2.2em]
\cdots \arrow[r] & V_{k+1} \arrow[r,"\partial_{k+1}"] \arrow[d,"f_{k+1}"'] & V_k \arrow[r,"\partial_k"] \arrow[d,"f_k"'] & V_{k-1} \arrow[r] \arrow[d,"f_{k-1}"'] & \cdots \\
\cdots \arrow[r] & W_{k+1} \arrow[r,"\xi_{k+1}"'] & W_k \arrow[r,"\xi_k"'] & W_{k-1} \arrow[r] & \cdots
\end{tikzcd}
\]

% block: chain-map-induced-homology page:25 type:ordinary-paragraph
\BT{In other words, a chain map is a graded linear map of degree 0 which commutes with the differentials $f\circ\partial=\xi\circ f$. Since $f(\partial c)=\xi f(c)$ for $c\in V_{k+1}$, $f$ maps the $k$-boundaries of $V$ to $k$-boundaries of $W$, and the same equation applied to $c\in\mathcal Z_k$ shows that $f$ maps $k$-cycles to $k$-cycles. It follows that the chain map $f$ induces a linear map $\bar f$ on homology, i.e., the map $\bar f_k:\mathcal H_k(V)\to\mathcal H_k(W)$ which takes the homology class $\bar v$ of some $v\in V_k$ to the homology class $\overline{f_kv}$ of $f_kv\in W_k$. This map is functorial. That is, if $(X,\zeta)$ is a third chain complex and $g$ is a chain map from $W$ to $X$, then $g\circ f$ is a chain map from $V$ to $X$ and $\overline{g\circ f}=\bar g\circ\bar f$.}{换言之, 链映射是次数为 0 且与微分交换的分次线性映射, 即 $f\circ\partial=\xi\circ f$. 由于对 $c\in V_{k+1}$ 有 $f(\partial c)=\xi f(c)$, $f$ 将 $V$ 的 $k$-边界链映为 $W$ 的 $k$-边界链; 对 $c\in\mathcal Z_k$ 应用同一等式, 可知 $f$ 将 $k$-闭链映为 $k$-闭链. 因而链映射 $f$ 在同调上诱导线性映射 $\bar f$, 即映射 $\bar f_k:\mathcal H_k(V)\to\mathcal H_k(W)$, 它将某个 $v\in V_k$ 的同调类 $\bar v$ 映为 $f_kv\in W_k$ 的同调类 $\overline{f_kv}$. 这一映射具有函子性. 也就是说, 若 $(X,\zeta)$ 是第三个链复形, 且 $g$ 是从 $W$ 到 $X$ 的链映射, 则 $g\circ f$ 是从 $V$ 到 $X$ 的链映射, 并且 $\overline{g\circ f}=\bar g\circ\bar f$.}

% block: chain-projection-definition page:25 type:ordinary-paragraph
\BT{A case we shall find of great interest later is when the complex $W$ is a subcomplex of $V$ and the maps $f_k$ are projections of $V_k$ onto $W_k$, i.e., $f_kv=v$ for $v\in W_k$. In this case we say that $f$ is a \emph{chain projection}.}{后面我们将特别关注这样一种情形: 复形 $W$ 是 $V$ 的子复形, 且映射 $f_k$ 是 $V_k$ 到 $W_k$ 上的投影, 即对 $v\in W_k$ 有 $f_kv=v$. 此时称 $f$ 为\emph{链投影}.}

% block: chain-projection-surjectivity page:25 type:proposition
\begin{Proposition}\label{prop:chain-projection-homology-surjection}
\BT{A chain projection from a complex to a subcomplex induces a surjection on homology.}{从复形到子复形的链投影在同调上诱导一个满射.}
\end{Proposition}

% block: chain-projection-surjectivity-proof page:25 type:proof
\begin{Proof}
\BT{The subcomplex inclusion maps $W_k\hookrightarrow V_k$ certainly form a chain map, $i$. When $f$ is a projection, $f\circ i$ is the identity map from $W$ to itself, and so $\bar f\circ\bar i=\overline{f\circ i}$ is the identity on homology and is, in particular, surjective. Therefore $\bar f$ must be surjective.}{子复形的包含映射 $W_k\hookrightarrow V_k$ 显然构成一个链映射 $i$. 当 $f$ 是投影时, $f\circ i$ 是从 $W$ 到自身的恒等映射, 因而 $\bar f\circ\bar i=\overline{f\circ i}$ 是同调上的恒等映射, 特别地, 它是满射. 因此 $\bar f$ 必为满射.}
\end{Proof}

% block: cochain-complexes-heading page:26 type:section
\subsection{\BT{Cochain complexes}{上链复形}}\label{sec:cochain-complexes}

% block: cochain-complex-definition page:26 type:ordinary-paragraph
\BT{A cochain complex is the same thing as a chain complex except with the indexing reversed: cochain differentials have degree $+1$ rather than $-1$. One typically writes the index of a cochain complex as a superscript rather than a subscript and often uses $d$ rather than $\partial$ to denote the differential:}{上链复形与链复形相同, 只是指标的方向相反: 上链微分的次数是 $+1$ 而非 $-1$. 上链复形的指标通常写作上标而非下标, 并且常用 $d$ 而非 $\partial$ 表示微分:}

% block: cochain-complex-sequence page:26 type:equation
\[
\cdots\to V^{k-1}\xrightarrow{d^{k-1}}V^k\xrightarrow{d^k}V^{k+1}\to\cdots.
\]

% block: cohomology-definition page:26 type:ordinary-paragraph
\BT{Of course, all the definitions and results above carry over in an obvious way. For example, the $k$th cohomology space $\mathcal H^k$ is defined as the quotient of the cocycles $\mathcal Z^k$ by the coboundaries $\mathcal B^k$.}{当然, 上述所有定义和结果都以显然的方式相应成立. 例如, 第 $k$ 个上同调空间 $\mathcal H^k$ 定义为闭上链 $\mathcal Z^k$ 模去上边界链 $\mathcal B^k$ 所得的商.}

% block: dual-cochain-complex page:26 type:ordinary-paragraph
\BT{Since any linear map $L:V\to W$ between vector spaces induces a dual map $L':W'\to V'$ between the dual spaces, any chain complex gives rise to a dual cochain complex. For example, from the simplicial chain complex associated to a simplicial complex, we obtain the simplicial cochain complex. Since the $k$-dimensional simplices form a basis for $C_k$, a $k$-cochain is an element of $C^k:=C'_k$, i.e., a function $\omega$ which assigns to each $k$-dimensional simplex a real number $\omega(f)$ (extended to chains by linearity). The differential $\partial^k:C^k\to C^{k+1}$ is defined by $\partial^k=\partial'_{k+1}$, i.e.,}{由于向量空间之间的任意线性映射 $L:V\to W$ 都在对偶空间之间诱导一个对偶映射 $L':W'\to V'$, 任意链复形都产生一个对偶上链复形. 例如, 从与单纯复形关联的单纯链复形出发, 我们得到单纯上链复形. 由于 $k$ 维单形构成 $C_k$ 的一组基, $k$-上链是 $C^k:=C'_k$ 的元素, 即一个函数 $\omega$, 它给每个 $k$ 维单形赋予一个实数 $\omega(f)$ (再由线性性扩张到链). 微分 $\partial^k:C^k\to C^{k+1}$ 由 $\partial^k=\partial'_{k+1}$ 定义, 即}

% block: dual-differential-evaluation page:26 type:equation
\[
(\partial^k\omega)g=\omega(\partial_{k+1}g),\qquad \omega\in C^k,\quad g\in C_{k+1}.
\]

% block: incidence-transpose page:26 type:ordinary-paragraph
\BT{This means that the matrix of the operator $\partial^k$ is the transpose of the oriented incidence matrix which represents $\partial_{k+1}$. For example, $\partial^0$ applied to a vertex is the sum of the edges incident on that vertex, each multiplied by $1$ or $-1$ according to whether the edge points into or out of the vertex. Of course, the information contained in the simplicial cochain complex}{这意味着算子 $\partial^k$ 的矩阵是表示 $\partial_{k+1}$ 的定向关联矩阵的转置. 例如, 将 $\partial^0$ 作用于一个顶点, 所得的是与该顶点关联的各条边之和, 每条边根据其指向该顶点还是背离该顶点而分别乘以 $1$ 或 $-1$. 当然, 单纯上链复形}

% block: simplicial-cochain-complex page:26 type:equation
\begin{equation}\label{eq:simplicial-cochain-complex}
0\to C^0\xrightarrow{\partial^0}C^1\xrightarrow{\partial^1}\cdots\xrightarrow{\partial^{n-1}}C^n\to0
\end{equation}

% block: cohomology-duality page:26 type:ordinary-paragraph
\BT{is essentially the same as that contained in the simplicial chain complex. In particular, it can be shown that the cohomology space $\mathcal H^k$ is naturally isomorphic to $\mathcal H'_k$, the dual of the corresponding homology space, and so its dimension is the $k$th Betti number.}{所包含的信息与单纯链复形所包含的信息本质上相同. 特别地, 可以证明上同调空间 $\mathcal H^k$ 自然同构于相应同调空间的对偶 $\mathcal H'_k$, 因而其维数是第 $k$ 个 Betti 数.}

% block: de-rham-heading page:26 type:section
\subsection{\BT{Example: The de Rham complex}{Example: de Rham 复形}}\label{sec:de-rham-complex}

% block: de-rham-introduction page:26 type:ordinary-paragraph
\BT{We will have a lot more to say about the de Rham complex}{我们还将更详细地讨论 de Rham 复形}

% block: differential-forms-complex page:26 type:equation
\begin{equation}\label{eq:differential-forms-de-rham-complex}
0\to\Lambda^0(\Omega)\xrightarrow{d^0}\Lambda^1(\Omega)\xrightarrow{d^1}\cdots\xrightarrow{d^{n-1}}\Lambda^n(\Omega)\to0
\end{equation}

% block: de-rham-background page:26 type:ordinary-paragraph
\BT{associated to any $n$-dimensional domain or smooth manifold $\Omega$. Here $\Lambda^k(\Omega)$ is the space of (smooth) differential $k$-forms on $\Omega$ and the operator $d$ is the exterior derivative, as introduced by Élie Cartan [\ref{bib:cartan-differential-expressions}]. We shall put off the introduction of the language of differential forms until Chapter~\ref{chap:exterior-calculus}. For now we just describe the case of $\Omega$ a domain in $\mathbb R^3$, for which the de Rham complex can be described in the terminology of vector calculus.}{它与任意 $n$ 维区域或光滑流形 $\Omega$ 关联. 这里 $\Lambda^k(\Omega)$ 是 $\Omega$ 上的 (光滑) 微分 $k$-形式空间, 算子 $d$ 是由 Élie Cartan [\ref{bib:cartan-differential-expressions}] 引入的外微分. 我们将微分形式语言的介绍留到 Chapter~\ref{chap:exterior-calculus}. 目前只描述 $\Omega$ 是 $\mathbb R^3$ 中区域的情形, 在这种情形下, de Rham 复形可以用向量分析的术语来描述.}

% block: vector-calculus-complex-intro page:26 type:ordinary-paragraph
\BT{The complex is then}{此时该复形为}

% block: vector-calculus-de-rham-complex page:26 type:equation
\begin{equation}\label{eq:vector-calculus-de-rham-complex}
0\to C^\infty(\Omega)\xrightarrow{\operatorname{grad}}C^\infty(\Omega;\mathbb R^3)\xrightarrow{\operatorname{curl}}C^\infty(\Omega;\mathbb R^3)\xrightarrow{\operatorname{div}}C^\infty(\Omega)\to0.
\end{equation}

% block: calculus-identities page:27 type:ordinary-paragraph
\BT{The calculus identities $\operatorname{curl}\operatorname{grad}\phi=0$ and $\operatorname{div}\operatorname{curl}u=0$ ensure that \eqref{eq:vector-calculus-de-rham-complex} is indeed a complex. Its cohomology spaces $\mathcal H^k_{\mathrm{dR}}$ comprise the de Rham cohomology of the domain.}{微积分恒等式 $\operatorname{curl}\operatorname{grad}\phi=0$ 和 $\operatorname{div}\operatorname{curl}u=0$ 保证 \eqref{eq:vector-calculus-de-rham-complex} 确实是一个复形. 它的上同调空间 $\mathcal H^k_{\mathrm{dR}}$ 构成该区域的 de Rham 上同调.}

% block: de-rham-map-zero-cochains page:27 type:ordinary-paragraph
\BT{If the domain $\Omega$ is triangulated by some simplicial complex $\mathcal S$, there is an interesting cochain map from the de Rham complex to the simplicial cochain complex called the \emph{de Rham map}. Namely, we associate to a function $\phi\in C^\infty(\Omega)$ the 0-cochain $F^0\phi$ determined by $F^0\phi(x)=\phi(x)$ for each $x\in\Delta_0(\mathcal S)$. As we shall see in Chapter~\ref{chap:exterior-calculus} a differential $k$-form may be integrated on an oriented $k$-dimensional manifold, and this is the simplest case of $k=0$. This is extended by linearity to arbitrary 0-chains, so $F^0\phi(\sum c_ix_i)=\sum c_i\phi(x_i)$. Similarly we associate to a vector field $\psi\in C^\infty(\Omega;\mathbb R^3)$ the 1-cochain whose value on any oriented edge $e\in\Delta_1(\mathcal S)$ is}{若区域 $\Omega$ 由某个单纯复形 $\mathcal S$ 给出三角剖分, 则存在一个从 de Rham 复形到单纯上链复形的有趣的上链映射, 称为 \emph{de Rham 映射}. 具体地, 我们给函数 $\phi\in C^\infty(\Omega)$ 配上一个 0-上链 $F^0\phi$, 它由对每个 $x\in\Delta_0(\mathcal S)$ 都有 $F^0\phi(x)=\phi(x)$ 确定. 正如我们将在 Chapter~\ref{chap:exterior-calculus} 看到的, 微分 $k$-形式可以在定向 $k$ 维流形上积分, 这里是最简单的 $k=0$ 情形. 由线性性将其扩张到任意 0-链, 便有 $F^0\phi(\sum c_ix_i)=\sum c_i\phi(x_i)$. 类似地, 我们给向量场 $\psi\in C^\infty(\Omega;\mathbb R^3)$ 配上一个 1-上链, 它在任意定向边 $e\in\Delta_1(\mathcal S)$ 上的值为}

% block: de-rham-edge-integral page:27 type:equation
\[
F^1\psi(e)=\int_e\psi\cdot t\,ds_e,\qquad e\in\Delta_1(\mathcal S),
\]

% block: de-rham-trace-forms page:27 type:ordinary-paragraph
\BT{where the $t$ denotes the unit vector tangent to the edge $e$ directed according to the orientation. The vector field $\psi$ may be viewed as a differential 1-form, and it has a trace on $e$ which is a differential 1-form on the one-dimensional manifold $e$ and so may be integrated over it. The line integral above is the expression of that integral in vector calculus terminology. A vector field in $\mathbb R^3$ may alternately be viewed as a differential 2-form, in which case the integral of its trace over a surface is defined and may be realized as}{其中 $t$ 表示沿边 $e$ 的定向取向的单位切向量. 向量场 $\psi$ 可以看作微分 1-形式, 它在 $e$ 上有迹, 这个迹是一维流形 $e$ 上的微分 1-形式, 因而可以在其上积分. 上述线积分就是该积分用向量分析术语表示的形式. $\mathbb R^3$ 中的向量场也可以看作微分 2-形式, 此时它的迹在曲面上的积分有定义, 并可表示为}

% block: de-rham-face-integral page:27 type:equation
\[
F^2\psi(f)=\int_f\psi\cdot n\,ds_f,\qquad f\in\Delta_2(\mathcal S),
\]

% block: de-rham-volume-form-intro page:27 type:ordinary-paragraph
\BT{where now $n$ is the unit normal to the face $f$, its direction determined by the orientation of $f$ according to the right-hand rule. Finally, a scalar function $\phi$ may be thought of as a 3-form, leading to}{这里 $n$ 是面 $f$ 的单位法向量, 其方向由 $f$ 的定向按照右手规则确定. 最后, 标量函数 $\phi$ 可以看作 3-形式, 从而得到}

% block: de-rham-volume-integral page:27 type:equation
\[
F^3\phi(T)=\int_T\phi\,dx
\]

% block: de-rham-connecting-map page:27 type:ordinary-paragraph
\BT{for a positively oriented tetrahedron $T\in\Delta_3(\mathcal S)$. We have thus defined the de Rham map connecting the de Rham complex \eqref{eq:vector-calculus-de-rham-complex} to the simplicial cochain complex \eqref{eq:simplicial-cochain-complex}:}{其中 $T\in\Delta_3(\mathcal S)$ 是正定向四面体. 这样, 我们就定义了将 de Rham 复形 \eqref{eq:vector-calculus-de-rham-complex} 与单纯上链复形 \eqref{eq:simplicial-cochain-complex} 连接起来的 de Rham 映射:}

% block: de-rham-commuting-diagram page:27 type:equation
\begin{equation}\label{eq:de-rham-commuting-map}
\begin{tikzcd}[column sep=1.1em,row sep=2.2em]
0 \arrow[r] & C^\infty(\Omega) \arrow[r,"\operatorname{grad}"] \arrow[d,"F^0"'] & C^\infty(\Omega;\mathbb R^3) \arrow[r,"\operatorname{curl}"] \arrow[d,"F^1"'] & C^\infty(\Omega;\mathbb R^3) \arrow[r,"\operatorname{div}"] \arrow[d,"F^2"'] & C^\infty(\Omega) \arrow[r] \arrow[d,"F^3"'] & 0 \\
0 \arrow[r] & C^0 \arrow[r,"\partial^0"'] & C^1 \arrow[r,"\partial^1"'] & C^2 \arrow[r,"\partial^2"'] & C^3 \arrow[r] & 0
\end{tikzcd}
\end{equation}

% block: de-rham-first-cell-verification page:27 type:ordinary-paragraph
\BT{The de Rham map is a cochain map, which is to say that this diagram commutes. For example, let us verify that $F^1\circ\operatorname{grad}=\partial^0\circ F^0$. We must show that, for an arbitrary smooth function $\phi$, the 1-cochains $F^1\operatorname{grad}\phi$ and $\partial^0F^0\phi$ coincide, i.e., that}{de Rham 映射是上链映射, 也就是说, 该图交换. 例如, 我们来验证 $F^1\circ\operatorname{grad}=\partial^0\circ F^0$. 必须证明, 对任意光滑函数 $\phi$, 1-上链 $F^1\operatorname{grad}\phi$ 与 $\partial^0F^0\phi$ 相同, 即}

% block: de-rham-first-cell-identity page:27 type:equation
\[
F^1\operatorname{grad}\phi([x_0,x_1])=\partial^0F^0\phi([x_0,x_1])
\]

% block: oriented-edge-unit-tangent pages:27,28 type:ordinary-paragraph
\BT{for any oriented edge $[x_0,x_1]$ in the complex. Letting $t$ denote the unit vector $(x_1-x_0)/|x_1-x_0|$, we have}{对该复形中的任意定向边 $[x_0,x_1]$ 都成立. 令 $t$ 表示单位向量 $(x_1-x_0)/|x_1-x_0|$, 则有}

% block: first-cell-calculation page:28 type:equation
\[
\begin{aligned}
F^1\operatorname{grad}\phi([x_0,x_1])
&=\int_{[x_0,x_1]}\operatorname{grad}\phi\cdot t
&&\text{\BT{(definition of $F^1$)}{($F^1$ 的定义)}}\\
&=\phi(x_1)-\phi(x_0)
&&\text{\BT{(fundamental theorem of calculus)}{(微积分基本定理)}}\\
&=F^0\phi([x_1]-[x_0])
&&\text{\BT{(definition of $F^0$)}{($F^0$ 的定义)}}\\
&=F^0\phi(\partial_1[x_0,x_1])
&&\text{\BT{(definition of $\partial_1$)}{($\partial_1$ 的定义)}}\\
&=\partial^0F^0\phi([x_0,x_1])
&&(\partial^0=\partial'_1).
\end{aligned}
\]

% block: remaining-cells page:28 type:ordinary-paragraph
\BT{The commutativity of the other two cells in the diagram is proven in the same way, using the Stokes theorem and the divergence theorem from vector calculus:}{图中另外两个方格的交换性可以用同样的方法证明, 其中用到向量分析中的 Stokes 定理和散度定理:}

% block: stokes-divergence-identities page:28 type:equation
\[
\int_f\operatorname{curl}\psi\cdot n=\int_{\partial f}\psi\cdot t,\qquad
\int_T\operatorname{div}\psi=\int_{\partial T}\psi\cdot n.
\]

% block: comparison-of-complexes page:28 type:ordinary-paragraph
\BT{The de Rham complex which constitutes the first row of the diagram \eqref{eq:de-rham-commuting-map} is very different from the simplicial cochain complex which is its second row. The former is composed of infinite dimensional function spaces connected by differential operators, while the latter consists of finite dimensional spaces with operators that are essentially combinatorial (multiplication by matrices whose entries are all $1$, $-1$, or $0$). However, the cohomology of the two complexes is identical. This statement is a form of de Rham's theorem.}{构成图 \eqref{eq:de-rham-commuting-map} 第一行的 de Rham 复形与构成第二行的单纯上链复形很不相同. 前者由通过微分算子连接的无限维函数空间组成, 而后者由有限维空间组成, 其算子本质上是组合性的 (乘以元素均为 $1$, $-1$ 或 $0$ 的矩阵). 然而, 两个复形的上同调相同. 这一陈述是 de Rham 定理的一种形式.}

% block: de-rham-isomorphism page:28 type:theorem
\begin{Theorem}{de Rham}\label{thm:de-rham-isomorphism}
\BT{The map on cohomology induced by the de Rham map is an isomorphism from the de Rham cohomology spaces onto the simplicial cohomology spaces.}{由 de Rham 映射在上同调上诱导的映射, 是从 de Rham 上同调空间到单纯上同调空间上的同构.}
\end{Theorem}

% block: singular-and-simplicial-cohomology page:28 type:remark
\begin{Remark}\label{rem:singular-simplicial-de-rham}
\BT{De Rham's theorem as just stated equates de Rham cohomology, which is constructed analytically using differential operators, with the more geometric simplicial cohomology, constructed from chains of simplices. Although simplicial homology is in some sense the simplest and most easily computed homology theory, topologists often prefer singular homology for its greater generality and theoretical flexibility. In particular, simplicial homology is defined only for spaces which can be triangulated (or are homeomorphic images of such spaces), while singular homology, being built from chains of singular simplices (arbitrary continuous images of simplices) is defined for any topological space. An early, quite elementary, result in most expositions of homology is that the two approaches give identical homology (or cohomology) spaces when both apply, that is, for triangulated spaces. De Rham's theorem, which is a deeper result, is then usually stated as giving an isomorphism between de Rham cohomology and singular cohomology on a smooth manifold.}{刚才陈述的 De Rham 定理将借助微分算子用分析方法构造的 de Rham 上同调, 与由单形的链构造的、更具几何性的单纯上同调等同起来. 虽然单纯同调在某种意义上是最简单且最易计算的同调理论, 拓扑学家却往往更倾向于奇异同调, 因为它具有更大的普遍性和理论灵活性. 特别地, 单纯同调只对能够三角剖分的空间 (或这类空间的同胚像) 有定义; 而奇异同调由奇异单形 (单形的任意连续像) 的链构造而成, 对任意拓扑空间都有定义. 在大多数同调论的论述中, 一个较早出现且相当初等的结果是: 当两种方法都适用时, 即对经过三角剖分的空间, 它们给出相同的同调 (或上同调) 空间. 随后通常将更深刻的 De Rham 定理表述为: 在光滑流形上, de Rham 上同调与奇异上同调之间存在同构.}
\end{Remark}

% block: de-rham-later-proof page:28 type:ordinary-paragraph
\BT{De Rham's theorem is a deep result. However, we shall see in Section~\ref{sec:whitney-forms-de-rham} that a version of it can be proven quite easily once we have assembled the tools of FEEC.}{De Rham 定理是一个深刻的结果. 然而, 我们将在 Section~\ref{sec:whitney-forms-de-rham} 看到, 一旦备齐 FEEC 的工具, 就可以相当容易地证明它的一个版本.}
